The Taylor series is a way to represent an analytic function near a point by an infinite power series (assuming the series converges to the function in that neighborhood). Truncating this series gives polynomial approximations of the function. Formally speaking, the Taylor series of a function is an infinite sum of the evaluations of the function's derivatives at a single point. Given function $f(X)$, its Taylor series centered at $a$ is expressed as follows:

$
f(a) + \dfrac{f'(a)}{1!}(X - a) + \dfrac{f''(a)}{2!}(X - a)^2 + \dfrac{f'''(a)}{3!}(X - a)^3 + \cdots = \sum\limits_{d=0}^{\infty}\dfrac{f^{(d)}(a)}{d!}(X - a)^d
$

Truncating after degree $D$ gives the $D$-th Taylor polynomial $T_D(X)$. Inside a neighborhood where the Taylor series represents $f$, the error tends to zero as $D\to\infty$, but it need not decrease at every successive degree or monotonically with the distance from $a$. For a real function with $D+1$ continuous derivatives on the interval joining $a$ and $X$, Taylor's theorem gives
\[
|f(X)-T_D(X)|\le\frac{\sup_{u}|f^{(D+1)}(u)|}{(D+1)!}|X-a|^{D+1},
\]
where the supremum is over that interval. Increasing the degree outside the series' convergence region gives no general accuracy guarantee.